\documentclass[11pt]{article}

\usepackage[a4paper,margin=1.05in]{geometry}
\usepackage{mathpazo}
\usepackage[T1]{fontenc}
\usepackage{amsmath,amssymb}
\usepackage{microtype}
\usepackage{booktabs}
\usepackage{tabularx}
\usepackage{enumitem}
\usepackage{parskip}
\usepackage{titlesec}
\usepackage{fancyhdr}
\usepackage{tcolorbox}
\usepackage[colorlinks=true,linkcolor=accent,urlcolor=accent,citecolor=accent]{hyperref}

\tcbuselibrary{skins,breakable}

\definecolor{accent}{HTML}{1F5C6B}
\definecolor{accentbg}{HTML}{DCE9EC}
\definecolor{softgray}{HTML}{6B6B6B}
\definecolor{warnrule}{HTML}{B8860B}
\definecolor{warnbg}{HTML}{FBF0D9}

\titleformat{\section}{\large\bfseries\color{accent}}{\thesection}{0.7em}{}
\titleformat{\subsection}{\normalsize\bfseries}{\thesubsection}{0.6em}{}
\setlist[itemize]{leftmargin=1.3em,itemsep=2pt,topsep=3pt}
\setlist[enumerate]{leftmargin=1.6em,itemsep=2pt,topsep=3pt}

\pagestyle{fancy}
\fancyhf{}
\renewcommand{\headrulewidth}{0.3pt}
\lhead{\small\color{softgray}The population generative model}
\rhead{\small\color{softgray}\thepage}

\newtcolorbox{notebox}{
  enhanced, breakable, colback=accentbg!55, colframe=accent,
  boxrule=0pt, leftrule=2.5pt, arc=1pt, left=10pt, right=10pt, top=8pt, bottom=8pt}

\newtcolorbox{assumpbox}{
  enhanced, breakable, colback=warnbg!60, colframe=warnrule,
  boxrule=0pt, leftrule=2.5pt, arc=1pt, left=10pt, right=10pt, top=8pt, bottom=8pt}

\newcommand{\reals}{\mathbb{R}}
\newcommand{\Ncal}{\mathcal{N}}
\newcommand{\Mset}{\mathcal{M}}
\newcommand{\Tr}{^{\!\top}}
\DeclareMathOperator{\Cov}{Cov}
\DeclareMathOperator{\Corr}{Corr}
\DeclareMathOperator{\diag}{diag}
\DeclareMathOperator{\tr}{tr}
\newcommand{\ctr}{\mathrm{ctr}}
\newcommand{\wid}{\mathrm{wid}}

%% Status of each term: does the fit estimate it, or not?
\newcommand{\INF}{\textcolor{accent}{\textbf{inferred}}}
\newcommand{\FIX}{\textcolor{softgray}{input}}
\newcommand{\DER}{\textcolor{softgray}{derived}}
\newcommand{\OBS}{\textcolor{softgray}{data}}
\newcommand{\STG}{\textcolor{softgray}{stage 1}}

\title{\bfseries The population generative model}
\author{}
\date{}

\begin{document}
\maketitle
\vspace{-2.5em}

\section{What the model does}

The model gives a distribution of parameters for a population of PDAC patients.
One draw from this distribution is one patient. The model then predicts what a
published study of these patients shows. The fit compares this prediction with
the numbers in the papers.

We never observe one patient. We observe a median and an interquartile range
from a cohort. The model must therefore run all the way from a population of
parameters to a cohort summary before it can compare anything.

\section{The terms}

Throughout, $\odot$ is an elementwise product and $\diag(v)$ is the diagonal
matrix with $v$ on its diagonal. A logarithm of a vector acts elementwise.

Each table gives the status of the term. \INF{} means that the fit estimates
it. \STG{} means that we compute it from the stage-1 output before the fit.
\FIX{} means that we set or build it before the fit by some other route. \DER{}
means that it is a function of the other terms. \OBS{} means that we read it
from a paper.

\subsection*{Sizes}
All sizes are constants.
\begin{center}\small
\begin{tabularx}{\textwidth}{@{}llX@{}}
\toprule
Term & Value & Meaning \\
\midrule
$P$ & 271 & The number of model parameters. \\
$M$ & 65 & The number of observables, over all scenarios. \\
$M_w$ & 49 & The number of observables whose source reports a usable scale. \\
$N$ & $2\times10^4$ & The number of patients that we simulate. Section \ref{sec:obs} sets it. \\
$C$ & 29 & The number of study cohorts. The pooled CD8 target splits into three; see section \ref{sec:repeat}. \\
$A$ & 118 & The number of anchor rows. It is $67$ centre rows plus $51$ width rows, because one observable now has three cohorts. \\
$n_c$ & 6--702 & The number of real patients in cohort $c$. \\
$B$ & $2\times10^3$ & Bootstrap replicates for the sampling covariance. \\
$B'$ & 200 & Independent clouds for the Monte Carlo covariance. \\
\bottomrule
\end{tabularx}
\end{center}

\noindent
Earlier versions of this note wrote $R$ for both the correlation matrix and the
number of anchor rows. The anchor count is now $A$. $R$ is only the matrix.

\subsection*{Quantities for one patient}
\begin{center}\small
\begin{tabularx}{\textwidth}{@{}lllX@{}}
\toprule
Term & Size & Status & Meaning \\
\midrule
$\theta_i$ & $P$ & \DER & The parameter vector of patient $i$. All values are positive. \\
$\vartheta_i=\log\theta_i$ & $P$ & \DER & The same vector on the log scale. \\
$z_i$ & $P$ & \FIX & The latent coordinates of patient $i$. \\
$x_i$ & $M$ & \DER & The observables of patient $i$, on the log scale. \\
\bottomrule
\end{tabularx}
\end{center}

\subsection*{Quantities for the population}
\begin{center}\small
\begin{tabularx}{\textwidth}{@{}lllX@{}}
\toprule
Term & Size & Status & Meaning \\
\midrule
$\mu$ & $P$ & \INF & The centre of the population, on the log scale. \\
$\omega$ & $P$ & \DER & The spread of each parameter between patients, on the log scale. \\
$s$ & 1 & \INF & The global width multiplier, on the log scale. It scales every assumed width at once. \\
$u$ & $P$ & \INF & The departure of each assumed width from the global level. It sums to zero. \\
$\Mset$ & 9 & \FIX & The set of parameters whose spread a source measured. These take no $s$ and no $u$. \\
$\omega_0$ & $P$ & \FIX & The prior centre for $\omega$. Section \ref{sec:omega} builds it in four layers. \\
$\Omega$ & $P\times P$ & \DER & The covariance of the population, on the log scale. \\
$R$ & $P\times P$ & \STG & The correlation matrix of the population. Section \ref{sec:R}. \\
$L_R$ & $P\times P$ & \FIX & The Cholesky factor of $R$, so that $R=L_R L_R\Tr$. \\
$D_\omega$ & $P\times P$ & \DER & $\diag(\omega)$, the matrix of marginal widths. \\
\bottomrule
\end{tabularx}
\end{center}

\subsection*{Quantities for the data}
\begin{center}\small
\begin{tabularx}{\textwidth}{@{}lllX@{}}
\toprule
Term & Size & Status & Meaning \\
\midrule
$g$ & map & \FIX & The model. It changes log parameters into log observables. \\
$Q^\ctr$ & $M$ & \DER & The median that the model predicts, per observable. \\
$Q^\wid$ & $M_w$ & \DER & The log interquartile range that the model predicts. \\
$y$ & $A$ & \OBS & The same two quantities, read from the papers. \\
$\eta$ & $C$ & \INF & An offset for each study. \\
$\mathcal{E}_c$ & rule & \FIX & The eligibility rule of cohort $c$, where the paper reports one. \\
$\delta$ & $M$ & \INF & A structural error in the centre, for each observable. \\
$\kappa$ & $M_w$ & \INF & A structural error in the width, for each observable with a width row. \\
$\kappa_0$ & 1 & \INF & The width error shared by every observable. \\
$\varepsilon$ & $M$ & \FIX & The held-out residual of the emulator, for each observable. \\
$\Sigma$ & $A\times A$ & \FIX & The covariance of the observed rows. The fit does not change it. \\
$\Sigma_{\text{MC}}$ & $A\times A$ & \FIX & The part of $\Sigma$ that comes from the finite cloud. It is common to every cohort. \\
\bottomrule
\end{tabularx}
\end{center}

\subsection*{Quantities that control the priors}
\begin{center}\small
\begin{tabularx}{\textwidth}{@{}lllX@{}}
\toprule
Term & Value & Status & Meaning \\
\midrule
$\mu_0$ & $P$ & \STG & The centre of the prior on $\mu$. Section \ref{sec:centre}. \\
$\Sigma_1$ & $P\times P$ & \STG & The covariance of the prior on $\mu$. The stage-1 covariance on the log scale. \\
$\tau_s$ & 0.5 & \FIX & The width of the prior on $s$. It trades directly against $\kappa_0$. \\
$\tau_\omega$ & 0.5 & \FIX & The width of the prior on $u$, on the log scale. \\
$\tau_{\omega,j}$ & --- & \FIX & The width of the prior on $\log\omega_j$ for $j\in\Mset$. It comes from the sample size $n_j$. \\
$\tau_\eta$ & --- & \INF & The scale of the study offsets. \\
$\tau_\delta$ & --- & \INF & The scale of the centre errors. \\
$\tau_\kappa$ & --- & \INF & The scale of the width errors. \\
$\phi$ & --- & \INF & All the quantities that the fit estimates. \\
\bottomrule
\end{tabularx}
\end{center}

\noindent
The full parameter vector is
\begin{equation*}
\phi=(\mu,\ u_{\text{raw}},\ s,\ \log\omega_\Mset,\
\eta_{\text{raw}},\ \delta_{\text{raw}},\ \kappa_{\text{raw}},\ \log\kappa_0,\
\log\tau_\eta,\ \log\tau_\delta,\ \log\tau_\kappa).
\end{equation*}

\section{How to make one patient}
\label{sec:patient}

\begin{enumerate}
\item Draw $z_i$ from a standard normal distribution. Do this one time only,
      before the fit.
\item Multiply by $L_R$. This step installs the correlations.
\item Multiply element $j$ of the result by $\omega_j$. This step sets the width
      of each parameter.
\item Add $\mu$.
\item The result is $\vartheta_i$.
\item Take the exponential of each element. The result is $\theta_i$.
\end{enumerate}

In one equation:
\begin{equation}
\vartheta_i \;=\; \mu \;+\; D_\omega\,L_R\,z_i,
\qquad z_i\sim\Ncal(0,I_P),
\qquad \theta_i=\exp(\vartheta_i).
\label{eq:patient}
\end{equation}

One cloud serves every cohort. Cohorts differ only by the eligibility rule
$\mathcal{E}_c$ of section \ref{sec:cohort}, which selects a subset of that one
cloud. This is deliberate twice over. It states that every study samples the same
population, which is the claim the model is making. It also makes the Monte Carlo
error common to every cohort, so section \ref{sec:obs} can model that error as one
object rather than as $C$ independent ones.

The covariance of the population follows directly:
\begin{equation}
\vartheta_i\sim\Ncal(\mu,\Omega),
\qquad
\Omega \;=\; D_\omega\,L_R L_R\Tr\,D_\omega \;=\; D_\omega\,R\,D_\omega .
\label{eq:Omega}
\end{equation}

This form separates the two halves of a covariance matrix. $R$ holds the
correlations and the fit does not change it. $\omega$ holds the marginal widths
and the fit estimates all $P$ of them. The construction is known as the
separation strategy.

Two properties follow, and both matter for interpretation.
\begin{align}
\sqrt{\Omega_{jj}} &= \omega_j
  && \text{$\omega_j$ is exactly the width of parameter $j$.}\\
\Corr(\Omega) &= R
  && \text{The correlations do not move when the fit moves $\omega$.}
\end{align}
A diagonal rescaling cannot change a correlation matrix, so the second property
holds at every value of $\omega$, not only at one point.

\begin{notebox}
\textbf{Why step 1 happens only one time.} We draw $z$ once and we keep it. The
map from $\phi$ to $Q(\phi)$ is then continuous and has no random noise. The
sampler needs this property to calculate a gradient. The same $N$ patients move
continuously as $\phi$ moves.

The map is continuous but it is not smooth. A quantile selects an order
statistic, and the selected patient changes when two patients swap rank. An
eligibility rule has the same defect at its threshold. Both give a piecewise
gradient, and only the few patients near each level carry it. We therefore use
a smooth quantile: weight patient $i$ by a kernel in its own rank, and take a
weighted quantile. \texttt{inference/importance.py} already supplies
\texttt{weighted\_quantile}. The bandwidth is a numerical choice and it must be
small against the width of the cloud.

The remaining cost is a fixed approximation error of order $N^{-1/2}$ in $Q$.
Section \ref{sec:obs} measures that error and adds it to $\Sigma$.
\end{notebox}

\section{How to predict the data}

\begin{enumerate}
\item Make $N$ patients with the steps in the previous section, one cloud per
      cohort.
\item Put each patient through the model $g$. This gives $N$ observable vectors
      $x_i=g(\vartheta_i)$.
\item Where the paper reports an eligibility rule, keep only the patients that
      satisfy it. Use a smooth weight at the threshold, not a hard indicator.
\item For each observable, take the smooth weighted quantiles at $0.25$, $0.50$
      and $0.75$.
\item Reduce them to two rows.
\end{enumerate}

The two rows are the centre and the log width:
\begin{equation}
Q^\ctr_m(\phi) \;=\; \widehat F^{-1}_{m,\phi}(0.5),
\qquad
Q^\wid_m(\phi) \;=\; \log\big\{\widehat F^{-1}_{m,\phi}(0.75)-\widehat
F^{-1}_{m,\phi}(0.25)\big\}.
\label{eq:Q}
\end{equation}
An observable whose source reports a centre only has a centre row and no width
row. There are $M=65$ observables with a centre and $M_w=49$ with a width. One
observable carries three cohorts, so the rows number $67$ and $51$, and $A=118$.
Section \ref{sec:repeat} explains the difference between an observable and a
row.

\begin{notebox}
\textbf{Why the rows are the centre and the log width.} An earlier version used
the three raw quantiles as three rows. Three changes follow from this one.

The sampling distribution of an interquartile range is skewed and is bounded
below by zero. At $n_c=6$ a normal error term on the raw quantiles is wrong. On
the log scale the same quantity is close to normal at every cohort size that we
have.

The width correction becomes additive. Section \ref{sec:obs} shows this. An
additive correction on one row makes the aliasing of section \ref{sec:widthid}
visible instead of implicit.

The count of width parameters falls from $65$ to $49$. Sixteen observables have
no width row, so their $\kappa_m$ was never identified and always returned its
prior. Those entries are now absent rather than reported.
\end{notebox}

There is no closed form for $F_{m,\phi}$. The model $g$ is a set of differential
equations and it is not linear, so the pushforward of a log-normal through it
belongs to no standard family. Step 4 is therefore a numerical calculation, not
a shortcut. In practice $g$ is a neural-network surrogate for the simulator,
which is what makes $C\times N$ evaluations per gradient step affordable.

\section{What the emulator must cover}
\label{sec:emulator}

The emulator sets a hard limit on what the fit can say, and the limit is a
property of the training pool and not of the sampler.

\textbf{Every parameter varies in the pool.} If the pool held a parameter fixed,
the emulator never saw that parameter move. The likelihood along that direction
is then not flat. It does not exist, and a $\theta$ that varies the parameter is
an extrapolation. The pool must therefore vary all $P$ widths that the fit
estimates.

\textbf{The pool must cover the flanks.} The rows in \eqref{eq:Q} read the cloud
at $0.25$ and $0.75$, not at the centre. A design that covers the centre of the
prior well and the tails poorly gives its worst predictions exactly where the
data enter.

\textbf{Measure $\varepsilon$ where the fit reads the emulator.} The held-out
residual of the training design is the wrong number to use. After the fit, draw
from $F(\theta\mid\hat\phi)$, run those draws through the full simulator, and
compare. If the error there exceeds the design error, the fit read the emulator
outside its region and the run is not valid.

\textbf{Report a resolution number per parameter.} With $P$ active inputs, a
direction can look flat because the emulator did not resolve it. That reads the
same as a flat likelihood and it is a different thing. A per-parameter
sensitivity of the emulator, measured on the pool, separates the two.

\section{The observation equation}
\label{sec:obs}

The model corrects its prediction before the comparison. One correction moves the
centre. The other changes the width.
\begin{align}
\widetilde Q^\ctr_m &\;=\; Q^\ctr_m(\phi) \;+\; \eta_{c(m)} \;+\; \delta_m ,
\label{eq:corrctr}\\
\widetilde Q^\wid_m &\;=\; Q^\wid_m(\phi) \;+\; \log\kappa_m .
\label{eq:corrwid}
\end{align}
Then
\begin{equation}
y \;=\; \widetilde Q(\phi) \;+\; e,
\qquad e\sim\Ncal(0,\Sigma).
\label{eq:obs}
\end{equation}

Each term answers one question about how the model can be wrong.

$\eta_c$ absorbs a measurement offset that applies to a whole study. $\delta_m$
absorbs a systematic error of the model in the centre of observable $m$.
$\kappa_m$ absorbs a systematic error of the model in the \emph{width} of
observable $m$, so $\kappa_m=1$ means no error and $\kappa_m=2$ means the model
produces half the spread that the study reports.

Without $\kappa$ the model has no way to be wrong about a width. Only $\omega$
could then close a gap in spread, and the fit would report a fault of the
simulator as diversity between patients. Section \ref{sec:widthid} explains what
separates the two.

\subsection{How we build \texorpdfstring{$\Sigma$}{Sigma}}
\label{sec:sigma}

$\Sigma$ has three parts. Each part comes from a different cause.
\begin{equation}
\Sigma \;=\; \Sigma_{\text{samp}} \;+\; \Sigma_{\text{MC}} \;+\;
\diag(\varepsilon^2).
\label{eq:Sigma}
\end{equation}

\textbf{$\Sigma_{\text{samp}}$: the cohort has a finite number of patients.} We
do not use an asymptotic formula. We resample the model's own cloud:
\begin{enumerate}
\item Take the cloud at the plug-in value $\phi_0$.
\item Draw $n_c$ patients from it with replacement.
\item Compute the rows of cohort $c$ from those $n_c$ patients, with the same
      quantile convention that the paper used.
\item Repeat $B$ times. Take the sample covariance across the rows of that
      cohort.
\end{enumerate}
This is exact at every $n_c$. It needs no density estimate and no bandwidth. It
carries the dependence between observables of the same cohort exactly, because
the cloud carries it. Rows of different cohorts get zero entries here, because
the real patients differ. $\Sigma_{\text{samp}}$ is therefore block diagonal by
cohort.

\begin{notebox}
\textbf{The formula, if you want one.} For a sample of size $n$ from any
continuous distribution $F$, the order statistic $X_{(k)}$ satisfies
$F(X_{(k)})\sim\mathrm{Beta}(k,\,n{+}1{-}k)$ exactly, at every $n$. Two levels
are jointly Dirichlet. So you can draw $u\sim\mathrm{Beta}(k,n{+}1{-}k)$, map
through the cloud's own quantile function $\widehat F^{-1}_{m,\phi_0}$, and read
the moments off the result. This is distribution free and it needs no
asymptotics.

The bootstrap above is the same calculation, and it also reproduces whatever
interpolation rule the paper applied. We prefer it for that reason.

The earlier version of this note used the Bahadur covariance, which needs a
density estimate $f_m$ at each quantile and a bandwidth $h=0.035$. That formula
is an asymptotic approximation to the quantity above. It is poor at $n_c=6$, and
some cohorts are that small.
\end{notebox}

\textbf{$\Sigma_{\text{MC}}$: the cloud is finite.} One cloud of $N$ patients
makes the prediction for all $C$ cohorts. Its error is therefore the same error
in every row, and not an independent error per cohort. We measure it directly.
Draw $B'$ independent clouds of size $N$ at $\phi_0$. Compute the full $A$-vector
$\widetilde Q$ for each cloud. Take the sample covariance across the $B'$ clouds.
The result is dense across cohorts. Keep it as a low-rank matrix plus a diagonal
when $B'<A$.

Report $\tr\Sigma_{\text{MC}}/\tr\Sigma_{\text{samp}}$. A reader must be able to
see whether the cloud mattered.

Reduce this term before you model it. Two actions do most of the work. Raise
$N$: the term falls as $N^{-1}$ and the emulator makes $N$ cheap. At
$N=2\times10^4$ against the largest cohort at $n_c=702$, it adds about $3.5\%$.
Stratify $z$ instead of drawing it independently: a scrambled Sobol sequence or
a stratified normal sample cuts the quantile variance at the same $N$.

\textbf{$\diag(\varepsilon^2)$: the emulator is not the simulator.} Section
\ref{sec:emulator} says how to measure $\varepsilon$ and where.

\begin{assumpbox}
\textbf{$\Sigma$ does not change during the fit.} The cloud that we resample
depends on $\phi$. We evaluate $\Sigma$ once, at a plug-in value $\phi_0$, and
then hold it fixed. The method is therefore an estimating equation, not a full
likelihood.

This is deliberate. If $\Sigma$ changes with $\phi$, then $\Sigma$ grows with the
population spread, and the term $-\tfrac12\log\det\Sigma$ rewards a narrow
population whatever the fit quality. A test showed the size of the error. With
cohorts of $20$, $40$ and $120$ patients, a live $\Sigma$ recovered $0.878$,
$0.941$ and $0.963$ of the true spread. A fixed $\Sigma$ recovered $1.002$,
$0.999$ and $1.001$.

The plug-in value is not critical. A change of $\phi_0$ by a factor of two moved
the estimate by $0.14$ to $0.32$ sampling standard deviations.

Take $\phi_0$ from the flat fit of section \ref{sec:flat}: the pair
$(\hat\mu_{\text{flat}},\,\omega_0)$. That fit costs about $1/N$ of this one, and
it is a better anchor than the prior mean.
\end{assumpbox}

\begin{assumpbox}
\textbf{The frozen cloud gives a bias, not a noise.} For one fixed $z_{1:N}$ the
Monte Carlo error is a fixed offset in $\widetilde Q$. It does not average away
over the $A$ rows. $\Sigma_{\text{MC}}$ is correct on average over clouds, which
is a hedge and not a cure.

Measure the residual bias directly. Run the whole fit three times, with three
different frozen clouds, and compare the posteriors. Section \ref{sec:checks}
lists this as a required check.
\end{assumpbox}

\section{The flat special case}
\label{sec:flat}

The flat problem asks for one parameter vector, not a population. It fits the
reported medians and it ignores the reported spreads. This model contains it,
and the reduction is worth stating, because we run the flat fit first.

\subsection{The reduction}

Set $\omega=0$. Then $D_\omega L_R z_i = 0$ for every $i$, all $N$ patients
collapse onto one point, and
\begin{equation}
Q^\ctr_m(\phi) \;=\; g_m(\mu).
\end{equation}
That is the flat forward map exactly.

The width rows do not survive the same limit. $Q^\wid_m=\log 0$ is not defined.
We remove them. The reduction is therefore a mask on the rows together with
$\omega=0$, and not a value of a parameter alone.

What remains is a small model:
\begin{equation}
y^\ctr_{m,c} \;=\; g_m(\mu) \;+\; \eta_c \;+\; \delta_m \;+\; e,
\qquad e\sim\Ncal(0,\Sigma),
\qquad \mu\sim\Ncal(\mu_0,\Sigma_1).
\label{eq:flat}
\end{equation}

\begin{center}\small
\begin{tabularx}{\textwidth}{@{}lX@{}}
\toprule
 & Terms \\
\midrule
Unchanged & $\mu$, $\mu_0$, $\Sigma_1$, $\eta$, $\delta$, $\tau_\eta$,
             $\tau_\delta$, $\varepsilon$ \\
Removed & $\omega$, $s$, $u$, $\omega_0$, $\Mset$, $R$, $L_R$, $\kappa$,
          $\kappa_0$, $\tau_\kappa$, $z_{1:N}$, $\Sigma_{\text{MC}}$ \\
Sections that go & \ref{sec:spread}, \ref{sec:widthid}, \ref{sec:omega},
                   \ref{sec:R}, \ref{sec:emit} \\
Rows & $67$ centre rows, against $118$ \\
Forward passes per gradient step & $1$, against $C\times N$ \\
\bottomrule
\end{tabularx}
\end{center}

Two parts of this document carry over word for word. Section \ref{sec:alias} on
the aliasing of $\eta$ and $\delta$ describes the map from observables to
cohorts, and that map does not change. Section \ref{sec:repeat} on the CD8 split
is the same, and it buys the same thing.

\subsection{Three things that the mask does not do}

\begin{assumpbox}
\textbf{The estimand is different, and one symbol hides it.} Flat $\mu$ is the
parameter vector that reproduces the reported medians. Population $\mu$ is the
centre of the population. For a nonlinear $g$ these differ, because
$g(\mu)\neq Q^\ctr(\mu,\omega)$ when $\omega\neq 0$. The difference is exactly
the Jensen gap of section \ref{sec:centre}.

So do not carry $\hat\mu_{\text{flat}}$ into the population model as a centre,
and do not read the population $\hat\mu$ as a flat calibration. Measure the gap
first, as section \ref{sec:centre} requires, and then translate.
\end{assumpbox}

\textbf{$\Sigma$ must come from another source.} Section \ref{sec:sigma}
resamples the model's cloud. The flat model has no cloud to resample. It has a
better source: the spread that the paper itself reports. The standard error of a
reported median follows from that spread and from $n_c$, with no plug-in and no
simulation.

Note the reversal. In the population model the interquartile range is signal. In
the flat model the same number is only noise calibration. The two fits therefore
use the same published table in two different ways, and a target that reports no
spread is weaker in both.

\textbf{The eligibility rule stops working.} A filter on one point does nothing.
The flat model therefore has no cohort device at all, and $\eta_c$ carries every
difference between studies on its own. The assumption box of section
\ref{sec:cohort} applies to the flat model without any qualification: the assay
story is asserted, and nothing in the flat fit can test it.

\subsection{Why we run it first}

\textbf{It is a test of everything except the population layer.} The prior on
$\mu$, the forward map, the corrections, the alias structure and the sampler all
run in \eqref{eq:flat}, at about $1/N$ of the cost. A fault in any of them is a
fault that the population fit would also have. Find it in the cheap model.

\textbf{It supplies the plug-in.} Section \ref{sec:sigma} holds $\Sigma$ at a
value $\phi_0$. Take $\phi_0=(\hat\mu_{\text{flat}},\,\omega_0)$.

\textbf{It measures the Jensen gap for free.} Compare $g(\hat\mu_{\text{flat}})$
against $Q^\ctr$ evaluated at the same centre with $\omega_0$. That is the number
that section \ref{sec:centre} asks for, and the flat fit is where to get it.

\eqref{eq:flat} is the population code with a mask, and it is not the neural
posterior estimator that answers the flat question elsewhere in this package.
The two routes take the same data to the same estimand. They do not share a
likelihood, a prior on the corrections, or a sampler. A disagreement between them
is informative, and it is not a defect by itself.

\section{The priors}

\subsection{The centre}
\label{sec:centre}

\begin{equation}
\mu \;\sim\; \Ncal\big(\mu_0,\;\Sigma_1\big).
\label{eq:muprior}
\end{equation}
$\mu_0$ and $\Sigma_1$ are the mean and the covariance of the stage-1 composite
prior, on the log scale. Where the stage-1 object is a copula, use its own
\texttt{log\_prob} as the prior term and skip the Gaussian approximation.

Earlier versions used $\Sigma_1=\tau_\mu^2 I_P$ with $\tau_\mu=0.5$. That gave a
parameter constrained to within $0.05$ and a parameter constrained to within
$1.5$ the same width, and it discarded the stage-1 correlations. Both are
available at no cost, so both are now used.

The width of \eqref{eq:muprior} measures how well we know the centre of the
population. It shrinks as data accumulate. It is a different quantity from
$\omega$, which is the variability between patients and does not shrink. To
confuse the two is the standard-error-versus-standard-deviation error, one level
up in the hierarchy.

\begin{assumpbox}
\textbf{Two conditions on $\mu_0$, and both are checkable.}

\emph{The Jensen gap.} Stage 1 estimated a parameter that reproduces a reported
summary. Under a nonlinear multivariate $g$, that parameter is not the parameter
of the median patient. So $g(\mu_0)$ and $Q^\ctr(\mu_0,\omega_0)$ differ.
Measure the difference per observable and compare it against $\tau_\delta$. If
it is small, record the number and continue. If it is not small, correct $\mu_0$
by one offline solve: find the shift $c$ with
$\mathrm{median}_i\{g(\mu_0+c+D_{\omega_0}L_R z_i)\}=g(\mu_0)$. That costs one
forward pass on the cloud.

\emph{Double use of the data.} \eqref{eq:muprior} is a prior on $\mu$ only if the
stage-1 sources are disjoint from the $C$ cohorts. If any publication feeds both,
$\mu_0$ counts that data twice and the posterior on $\mu$ comes out too narrow.
The source lists make this checkable, and it has not yet been checked.
\end{assumpbox}

\subsection{The spread}
\label{sec:spread}

The widths split into two sets. A parameter is in $\Mset$ when a source measured
its spread across biological units, in the sense of section \ref{sec:omega}. At
present $|\Mset|=9$.

For a measured parameter, $j\in\Mset$:
\begin{equation}
\log\omega_j \;\sim\; \Ncal\big(\log\omega_{0j},\;\tau_{\omega,j}^2\big),
\qquad
\tau_{\omega,j}^2 \;=\; \frac{1}{2(n_j-1)+\tau^{-2}} .
\label{eq:omegameas}
\end{equation}
The width is the posterior width that the shrinkage \eqref{eq:shrink} already
implies. A measured parameter takes no global multiplier.

For an assumed parameter, $j\notin\Mset$:
\begin{equation}
\log\omega_j \;=\; \log\omega_{0j} \;+\; s \;+\; u_j,
\qquad
s\sim\Ncal(0,\tau_s^2),
\qquad
u_j\sim\Ncal(0,\tau_\omega^2),
\qquad
\sum_{j\notin\Mset} u_j = 0 .
\label{eq:omegaprior}
\end{equation}
We impose the constraint by centring: $u = u_{\text{raw}} - \bar u_{\text{raw}}$.

The split is the point. $s$ is one number that moves every assumed width
together. $u$ holds the pattern and has $P-|\Mset|-1$ free values. Section
\ref{sec:widthid} needs the two to be separate quantities, because $s$ is the one
that competes with $\kappa_0$ and $u$ is not.

The value $\tau_\omega$ controls how far the data can move one width away from
its anchor. On the log scale, $\tau_\omega=0.5$ gives a central $95\%$ interval
of about $[0.37,\,2.72]$ for each parameter.

This choice is not harmless. The data give at most $M_w=49$ constraints on
$\Omega$ (section \ref{sec:whatdataconstrain}), against $P=271$ widths. Most
parameters therefore return their prior. A large $\tau_\omega$ reports a wide
population that no study measured.

\subsection{The corrections}

\begin{align}
\eta &= \tau_\eta\,\eta_{\text{raw}}, &
\eta_{\text{raw}} &\sim\Ncal(0,I_C), &
\log\tau_\eta &\sim\Ncal(\log 0.15,\,1),\\
\delta &= \tau_\delta\,\delta_{\text{raw}}, &
\delta_{\text{raw}} &\sim\Ncal(0,I_M), &
\log\tau_\delta &\sim\Ncal(\log 0.15,\,1),\\
\log\kappa &= \log\kappa_0 + \tau_\kappa\,\kappa_{\text{raw}}, &
\kappa_{\text{raw}} &\sim\Ncal(0,I_{M_w}), &
\log\tau_\kappa &\sim\Ncal(\log 0.15,\,1),
\end{align}
with $\log\kappa_0\sim\Ncal(0,\,\tau_s^2)$. These are non-centred forms. The
centre $0.15$ matches the floor of the translation-relevance rubric used at
stage 1.

The prior on $\log\kappa_0$ takes the same width as the prior on $s$. The two
compete for one direction, and giving them the same width states that we have no
preference between them before the data arrive.

\subsection{The cohort terms}
\label{sec:cohort}

Studies do not recruit the same patients. A cohort of resectable patients and a
cohort of metastatic patients are not two samples from one distribution, so the
claim that one $F$ generates all $C$ cohorts is a real assumption and not a
formality. This subsection says what the model does about it, and what it
deliberately does not do.

All but one observable belong to exactly one cohort. Section \ref{sec:alias}
shows that a \emph{free} per-cohort term on the observation is therefore aliased
with $\delta$, exactly, for those observables. Only two things escape that: a
term that costs no parameter, or an observable that appears in more than one
cohort. Section \ref{sec:repeat} is the second case.

\subsubsection*{The eligibility rule}

Cohort $c$ recruited patients that satisfy some criterion. Write $\mathcal{E}_c$
for the region of observable space that the criterion defines. The cohort is then
the same population under a condition, and we read its quantiles from the part of
the cloud that meets the condition:
\begin{equation}
\widehat F_{m,\phi,c}(t) \;=\;
\frac{\sum_i w^{(c)}_i\,\mathbf 1\{x_{im}\le t\}}{\sum_i w^{(c)}_i},
\qquad
w^{(c)}_i \;=\; \mathbf 1\{x_i\in\mathcal{E}_c\}.
\label{eq:elig}
\end{equation}
Then $Q^\ctr_{m,c}$ and $Q^\wid_{m,c}$ come off that weighted distribution
function. At the boundary, use a smooth weight
$w^{(c)}_i=\sigma\{(x_{ik}-a)/h\}\,\sigma\{(b-x_{ik})/h\}$ rather than a hard
indicator, for the reason that the note in section \ref{sec:patient} gives. The
bounds $a$ and $b$ come from the methods section of the paper.

The device has three properties, and the second one is the reason to want it.

\textbf{It costs no parameter.} The bounds are read, not estimated. There is
nothing for $\delta_m$ to trade against, so the aliasing of section
\ref{sec:alias} does not engage at all.

\textbf{It moves the width, and not only the centre.} $\eta_c$ shifts the median
of a cohort and can do nothing else. A truncation changes the median \emph{and}
the interquartile range, for the observable that carries the criterion and,
through the population correlations, for the other observables of that cohort. It
is the only device in this model that can make one cohort narrower than another.
Section \ref{sec:widthid} is entirely about a shortfall in width, so this matters
more than the shift does.

\textbf{It is the literal description of what the study did.} The population does
not change. The study observed a part of it and reported which part.

\begin{assumpbox}
\textbf{Three limits, and the first one is not yet measured.}

\emph{Most criteria cannot be expressed.} Equation \eqref{eq:elig} needs a
criterion that is a region of \emph{observable} space. Many are not.

\begin{center}\small
\begin{tabularx}{0.93\textwidth}{@{}Xl@{}}
\toprule
Criterion & Usable \\
\midrule
Tumour diameter between two values & yes, if burden is an observable \\
Resectable against metastatic & in part, if it maps to a burden threshold \\
Performance status, age & no, the model does not represent them \\
Treatment naive & no, that is a different scenario \\
Mutation status & no, that is a covariate on parameters \\
Biopsy accessible & no, there is no model correlate \\
\bottomrule
\end{tabularx}
\end{center}

Nobody has counted how many of the $29$ cohorts state a criterion in the first
group. That count is the measure of how much this device can carry, and it is
unknown. Make the count before relying on it. If it is small, then $\eta_c$
carries almost every difference between studies, and assumption (vi) carries the
model.

\emph{An inferred criterion is not free.} A study that used resection specimens
selected on resectability, whether or not it says so. A filter applied on that
basis is an assumption in the costume of a reported fact. Record which rules were
read and which were inferred, and treat the inferred ones as assumptions.

\emph{The cross-observable part inherits $R$.} A filter on observable $k$ moves
observable $m$ only through their dependence across patients, which comes from
$g$ and from $R$. Section \ref{sec:R} says $R$ is assumed and close to the
identity. So the effect on the filtered observable is sound, and the effect on
the other observables of the cohort is real but its size is assumed.
\end{assumpbox}

\textbf{Meta-regression, once the attribute table exists.} Replace the free
$\eta_c$ by $\eta_c = Z_c\beta$, where $Z_c$ holds cohort attributes: stage, line
of therapy, sample type, assay platform, region. The free count then falls from
$C$ to $\dim\beta$, and cohorts that share attributes share the effect. The
coefficients pool across all $C$ cohorts, which is what makes them estimable. We
do not have the table yet, so $\eta_c$ stays free for now.

The same device can be moved one level down, onto the parameters:
$\mu_c=\mu+\Gamma z_c$, with $\Gamma$ pooled across cohorts and each column tied
to a declared mechanism, such as stage moving a growth rate. That is the ordinary
covariate model of an NLME analysis. It is the correct upgrade path, and it is
also gated on the attribute table. Declare the links one at a time and from
biology. Do not fit a basis.

\begin{assumpbox}
\textbf{What $\eta_c$ asserts, and what we cannot test.} An additive offset on
every log readout of one study is an \emph{assay calibration} story: this
platform reads high by a constant factor. A cohort-level shift in the parameters
would be a \emph{biology} story: this study recruited different patients. The two
predict the same thing for a cohort with one observable, and most cohorts have
about two. These data cannot separate them.

The model asserts the assay story by construction. That is a choice, it is not a
finding, and any claim about $\mu$ inherits it. Section \ref{sec:data} lists the
data that would let the two be told apart.
\end{assumpbox}

\begin{notebox}
\textbf{Why there is no cohort random effect on the parameters.} An earlier
version carried $\mu_c=\mu+Sb_c$, with $b_c\sim\Ncal(0,\psi^2I_q)$ and $S$ a
fixed $P\times q$ loading. It was removed.

The construct itself is ordinary: a cohort random effect is the second level of
any NLME model. The failure was $S$. Its columns must say which directions in
parameter space a cohort may move along, and no evidence specifies them. Choosing
them by hand encodes a guess as a structural constraint. Choosing $S=I_P$ adds a
second variance component that competes with $\omega$ and that these data cannot
resolve.

The decisive objection is one of consistency. Section \ref{sec:whatdataconstrain}
deletes the basis matrix $W$, and the reason given is that a low-rank loading in
parameter space, whose columns we chose rather than measured, is not a biological
statement. $S$ is the same kind of object and it fails the same test. It was
introduced three sections after its twin was rejected.

The identification was thin in any case. A cohort holds about two observables, so
two rows had to separate a $q$-dimensional mechanistic shift from a free
$\delta_m$ per observable plus a shared $\eta_c$. Removing $b_c$, $S$, $\psi$ and
$q$ costs nothing that was ever identified. The covariate model above recovers
the same intent with coefficients that pool, and the eligibility rule recovers
part of it for free.
\end{notebox}

\subsection{The repeated endpoint}
\label{sec:repeat}

One endpoint does appear more than once. The pooled CD8 target combines three
cohorts into one row. Split it back out. This is the single most valuable change
in this document, and it is extraction work rather than modelling work.

Write $m^\ast$ for that observable and $c_1,c_2,c_3$ for its cohorts. The three
centre rows are
\begin{equation}
\widetilde Q^\ctr_{m^\ast,c_k} \;=\; Q^\ctr_{m^\ast}(\phi) \;+\; \eta_{c_k} \;+\;
\delta_{m^\ast},
\qquad k=1,2,3 .
\end{equation}
The three rows share one $\delta_{m^\ast}$ and differ only in $\eta$. Four things
follow.

\textbf{Two contrasts of $\eta$ become identified.} The differences
$\eta_{c_1}-\eta_{c_2}$ and $\eta_{c_1}-\eta_{c_3}$ are read directly from the
data. Only the mean of the three stays aliased with $\delta_{m^\ast}$. This is
the first identified study effect in the model.

\textbf{$\tau_\eta$ gets one direct measurement.} Two contrasts are a weak
estimate of a scale, and they are the only estimate available. $\tau_\eta$ then
sets how much of every \emph{other} residual in the model goes to a study effect
rather than to a model error. The value of these two contrasts therefore reaches
far past their own three rows.

\textbf{The pooled width was too large.} A pool of three cohorts with different
centres is a mixture. Its interquartile range holds the between-patient spread
\emph{and} the spread of the three centres. The pooled row therefore fed
$\omega$ with a width that no single cohort has. The shortfall figure of section
\ref{sec:widthid}, $0.62$ of the observed interquartile range, is biased by this
in the direction that makes the simulator look narrower than it is. Recompute it
after the split.

\textbf{The split is a variance-components decomposition.} Within a sub-cohort
the width is between-patient variability. Across the three sub-cohorts the
medians differ by a study effect plus a sampling error. This is the one place in
the whole dataset where both quantities are visible for the same endpoint. It is
also a direct test of $\Sigma_{\text{samp}}$: the disagreement of the three
medians must match what the bootstrap of section \ref{sec:sigma} predicts.

Two conditions apply. The sources must report a per-cohort $n$, median and
interquartile range, not only the pool. And the three sub-cohorts are small, so
the asymptotic quantile covariance would have been a poor description of them.
The bootstrap form of $\Sigma_{\text{samp}}$ is what makes the split affordable.

If the three sub-cohorts report a centre only, then $A=116$ rather than $118$,
and the width argument above still holds for the pooled row that they replace.

Search the remaining sources for the same pattern. Every further repeated
endpoint moves section \ref{sec:alias} further from a prior structure and closer
to a measurement.

One further question is worth asking of these three cohorts alone. Do any of them
report a second observable? If they do, then the three cohorts hold a repeated
endpoint \emph{and} more than one readout, which is the only configuration in the
whole set that could separate an assay effect from a recruitment effect. That is
the condition in assumption (vi), and this is the one place where it might
already be satisfied. Check it before accepting that assumption as unavoidable.

\subsection{What separates a narrow model from a narrow population}
\label{sec:widthid}

Both $\kappa$ and $\omega$ can widen a predicted interquartile range. The rows in
\eqref{eq:corrwid} make the competition exact and one dimensional, so we can
state it rather than argue about it.

To first order the variance of observable $m$ across patients is
\begin{equation}
\mathrm{Var}[x_m] \;\approx\; \sum_{j,k} J_{mj}J_{mk}R_{jk}\,\omega_j\omega_k,
\qquad J_m=\partial g_m/\partial\vartheta .
\end{equation}
Substitute \eqref{eq:omegaprior}. Every assumed width carries the same factor
$e^{s}$. Write $\lambda_m$ for the share of that variance that comes from the
measured set $\Mset$. Then
\begin{equation}
Q^\wid_m \;\approx\; \text{const}_m \;+\; (1-\lambda_m)\,s,
\qquad\text{so}\qquad
\widetilde Q^\wid_m \;\approx\; \text{const}_m \;+\; (1-\lambda_m)\,s \;+\;
\log\kappa_0 \;+\; \tau_\kappa\kappa_{\text{raw},m}.
\label{eq:ridge}
\end{equation}

Three statements follow, and each one is checkable.

\textbf{If $\lambda_m=0$ for every observable, then $s$ and $\log\kappa_0$ are
aliased exactly, to first order.} Only their sum appears. The prior then decides
the split, and nothing else does. This is the honest reading of the current
model.

\textbf{The break is $\lambda_m$, and it must vary across observables.} A
measured width takes no $s$, by \eqref{eq:omegameas}. An observable driven mostly
by $\Mset$ therefore reads $\log\kappa_0$ almost directly, and the difference
between a high-$\lambda$ observable and a low-$\lambda$ one separates the two.
Compute $\lambda_m$ for all $M_w$ observables and report the spread before you
quote the split. With $|\Mset|=9$ the values may all be near zero, in which case
the split is a prior statement and must be reported as one.

\textbf{The earlier argument was wrong in a way that mattered.} It compared
$\kappa_0$ against one $\omega_j$ and concluded that a uniform shortfall is cheap
under $\kappa_0$, because no single parameter feeds every observable. A uniform
rescaling of \emph{all} widths does feed every observable, exactly. What
separated the two was the prior mass: a common shift of $P$ independent
$\log\omega_j$ costs $P(\log c)^2/2\tau_\omega^2$, which grows with $P$. The
split therefore depended on how many inert parameters stayed in the vector.
Removing $200$ unused parameters would have moved it. The form
\eqref{eq:omegaprior} removes that dependence, because the global direction is
now one parameter with one prior width $\tau_s$.

The current model produces about $0.62$ of the observed interquartile range at
the median observable. Under \eqref{eq:ridge} that uniform shortfall lands on the
ridge $s+\log\kappa_0$. Report the joint posterior of $(s,\log\kappa_0)$ as a two
dimensional figure. A ridge is a result. A marginal that hides a ridge is not.

Recompute the $0.62$ after the CD8 split of section \ref{sec:repeat}. A pooled
cohort reports the width of a mixture, which is wider than the width of any
component. Any pooled row therefore pushes this figure down, and it does so for
a reason that has nothing to do with the simulator.

A third action attacks the same problem from the data side. With three quantile
levels, $\kappa_m$ copies any width change exactly, because it stretches the
prediction about its own median. With five levels, or a reported range, or a
reported mean beside a reported median, a change in $\omega$ changes the
\emph{shape} of the pushforward and $\kappa$ cannot copy it. Every extra level
that a source gives is a direct attack on this aliasing, and it costs nothing but
extraction work.

Per observable the split is still not identified. Each observable with a width
row has one centre row and one width row. $\delta_m$ takes the centre and
$\kappa_m$ takes the width, which leaves nothing. So $\tau_\delta$ and
$\tau_\kappa$ carry the separation, exactly as in section \ref{sec:alias}.

\section{How we build \texorpdfstring{$\omega_0$}{omega0}}
\label{sec:omega}

The vector $\omega_0$ is the prior centre for the marginal widths in
\eqref{eq:omegameas} and \eqref{eq:omegaprior}. We build it in four layers. Each
layer replaces the layer before it.

\begin{enumerate}[label=\textbf{L\arabic*.}]
\item \textbf{Default.} $\omega_j=0.35$ for every parameter. This is a
      coefficient of variation of about $36\%$. A parameter gets this value when
      we know nothing more specific.
\item \textbf{Class.} A geometric packing limit gets $0.15$. A material constant
      gets $0.10$. The spread that we observe for these classes is mostly assay
      noise. To call it $\omega$ would overstate the real diversity.
\item \textbf{Explicit.} One value for one parameter, when we hold a specific
      belief about it.
\item \textbf{Data.} A submodel target can report a spread across biological
      units. We then combine that spread with the value from L1 to L3. We do not
      replace it. A parameter that reaches L4 enters the set $\Mset$.
\end{enumerate}

For L4, the sampling standard deviation of a log spread from $n$ units is about
$\{2(n-1)\}^{-1/2}$. The data therefore carry precision $2(n-1)$ against a
class-prior precision $\tau^{-2}$. This gives
\begin{equation}
\log\omega_{0j} \;=\;
\frac{2(n_j-1)\,\log\hat s_j \;+\; \tau^{-2}\,\log\omega^{(0)}_j}
     {2(n_j-1)\;+\;\tau^{-2}},
\qquad \tau=0.5 .
\label{eq:shrink}
\end{equation}
In this equation, $\hat s_j$ is the spread that the source reports, $n_j$ is the
number of biological units behind it, and $\omega^{(0)}_j$ is the value from L1
to L3. The data become stronger than the prior at about $n=3$. The precision in
the denominator is also the precision that \eqref{eq:omegameas} uses as the prior
width for $j\in\Mset$.

A parameter with no value of $n_j$ keeps $\omega^{(0)}_j$. A spread with no
sample size carries no information.

\begin{assumpbox}
\textbf{What L4 admits.} Equation \eqref{eq:shrink} applies only when the source
measures variation across biological units, such as patients, donors or animals.
It does not apply to technical replicates, to translation between experimental
contexts, or to an assumed value. Those are measurement or epistemic quantities
and not variability between patients.

This gate matters more than it appears. The submodel population block covers
$178$ parameters, but for $145$ of them the reported spread is numerically equal
to the epistemic posterior width. Without the gate, the model would take those
$145$ epistemic widths as between-patient spread. At present the gate admits $9$
of $271$ parameters. The other $262$ come from L1 to L3.

The gate now does a second job. $\Mset$ is exactly the set that L4 admits, and
section \ref{sec:widthid} needs $\Mset$ to be non-empty to identify $\kappa_0$
at all. Every parameter that the gate admits therefore buys identification as
well as a width. Widening the gate to buy identification would defeat it.
\end{assumpbox}

\section{Where \texorpdfstring{$R$}{R} comes from}
\label{sec:R}

$R$ is a separate input to the model. Nothing in the fit changes it, and you can
replace it without touching anything else. At present it is the correlation
matrix of the stage-1 composite prior,
\begin{equation}
R \;=\; \Corr\big(\Cov_\pi[\vartheta]\big).
\label{eq:R}
\end{equation}
The stage-1 prior joins submodel posteriors through a Gaussian copula, so
\eqref{eq:R} carries whatever structure that fit produced. The copula covers
$65$ of the $271$ parameters. The remaining $206$ parameters come from
independent marginals, so $R$ states independence for every pair that involves
one of them.

\begin{assumpbox}
\textbf{What \eqref{eq:R} does and does not give.} Two measured facts qualify it.

The matrix is close to the identity. Its copula covers $65$ of $271$
parameters, and only $26$ off-diagonal entries exceed $0.5$ in absolute value,
out of $36{,}585$. For most parameter pairs, \eqref{eq:R} states independence.

Where the matrix is not close to the identity, the entry is epistemic. The
largest, at $-0.994$, joins two rates that one target constrains only through
their sum. That number describes what the stage-1 data could not separate. It
does not describe how patients differ. A second group of strong entries comes
from declared hierarchical parameter groups; those encode a real claim that the
members are the same kind of quantity, but their magnitudes are still epistemic.

$R$ is therefore best read as a placeholder with the right interface, not as a
measured quantity.
\end{assumpbox}

Better candidates for $R$ exist and slot into the same place:
\begin{itemize}
\item a block matrix built from the declared parameter groups, with a within-group
      correlation elicited rather than read from stage 1;
\item structural constraints, such as rates that partition a conserved total;
\item the identity, if the conclusion is that nothing is known;
\item a shrunk form $(1-\lambda)I + \lambda\,\Corr(\Cov_\pi[\vartheta])$.
\end{itemize}
Because $R$ enters only through \eqref{eq:patient}, a comparison across these
choices is a sensitivity analysis and not a redesign. Section \ref{sec:checks}
lists that comparison, and it has not yet been run.

\begin{notebox}
\textbf{A correlation is not a constraint.} Some parameters partition a conserved
total, so their values must satisfy an identity. $R$ cannot state an identity. A
correlation of $-1$ makes two parameters move together on a line, but the line is
not the constraint surface, and a lognormal draw does not sit on it. Only a
reparameterisation can enforce such a constraint. Section \ref{sec:emit} filters
the draws that break one, which detects the problem and does not fix it.
\end{notebox}

\section{What the data constrain}
\label{sec:whatdataconstrain}

The data constrain $\Omega$ only through the width of an observable
distribution. To first order the spread of observable $m$ across patients is
$J_m\Omega J_m\Tr$. Each width row therefore gives one linear function of
$\Omega$.

Published summaries are marginal, one observable at a time. The cross terms
$J_j\Omega J_l\Tr$ are never observed. Of the $65$ targets, $49$ report a usable
scale. There are therefore at most $49$ constraints on $\Omega$, and $\kappa$
takes some of them. There are $51$ width rows, because the CD8 target of section
\ref{sec:repeat} contributes three. Rows and constraints are different counts:
those three rows measure the same function of $\Omega$ three times, which buys
precision and not a new direction.

Two consequences follow.

First, no part of $R$ can be estimated from these data. Its entries are exactly
the cross terms that a marginal summary never reports. $R$ must come from
outside the fit, which is why section \ref{sec:R} treats it as an input.

Second, most of the $271$ widths are not informed either. Along those
directions the posterior returns the prior \eqref{eq:omegaprior}. The reported
population is then a statement of $\omega_0$, not a measurement.

\begin{notebox}
\textbf{Design note: there is no basis matrix.} An earlier version replaced
\eqref{eq:patient} by $\vartheta_i=\mu+W(\sigma_u\odot z_i)$. Here $W=L\widetilde
V$ came from an eigendecomposition of the observation information in the metric
of the prior covariance, and the fit estimated $K=15$ direction multipliers
instead of $271$ widths. The construction is known as the likelihood-informed
subspace, or the active subspace.

The rotation and the truncation were removed. The correlations were kept, and
they are now $R$.

Three reasons. First, $\sigma_u$ was not interpretable: it scaled directions, so
no single entry was the width of any parameter. Second, the correlations did not
stay fixed. Because $\Omega=\Gamma+\sum_{k\le K}(\sigma_{u,k}^2-1)w_kw_k\Tr$, the
fit moved $\Corr(\Omega)$ by a rank-$K$ update, and that update lay along the
directions that the \emph{measurements} resolve. A correlation that appears
because a combination happened to be observable is not a biological statement.
Third, the truncation to $K$ directions limited how much model error the fit
could report as patient diversity. That is protection by restriction and not by
correctness; section \ref{sec:widthid} now names the correct fix.

The present form keeps what the rotation was supposed to deliver, namely the
prior correlations, and states them as an explicit and replaceable object.
\end{notebox}

\section{The corrections are aliased}
\label{sec:alias}

Sixty-four of the $65$ observables belong to exactly one cohort. Only the CD8
target of section \ref{sec:repeat} appears in more than one. That fact has a
consequence that the notation hides, and it applies to all but one observable.

Take a cohort $c$ that holds observables $m_1,\dots,m_k$. Equation
\eqref{eq:corrctr} uses only the sum $\eta_c+\delta_{m_i}$. Add a constant to
$\eta_c$, subtract the same constant from every $\delta_{m_i}$ in that cohort,
and the likelihood does not change. $\eta$ and $\delta$ are therefore aliased,
exactly, and only their priors separate them.

The eligibility rule of section \ref{sec:cohort} is the one cohort device that
escapes this. It enters before $g$ and it costs no parameter, so there is nothing
for $\delta$ to trade against. A cohort effect that carries a free parameter does
not escape, whichever level it acts on, which is why section \ref{sec:cohort}
removed the parameter version rather than tuning it.

\begin{assumpbox}
\textbf{What $\eta$ is, and is not.} In an ordinary meta-analysis, a study
effect is identified because the same endpoint appears in several studies. One
endpoint does that here, and it gives two contrasts. For the other $26$ cohorts
$\eta_c$ is not an identified study effect.

Read those as a correlation structure on the corrections: the model states that
the errors within one study share a common part. That claim has content, and it
is a reasonable claim, but it is a prior structure and not a measurement.

What changes with the CD8 split is the \emph{scale}. $\tau_\eta$ is now measured
rather than asserted, and $\tau_\eta$ decides how much of every other cohort's
residual is a study effect. Two contrasts are a thin basis for that, so report
the posterior for $\tau_\eta$ beside any statement that depends on it. Section
\ref{sec:cohort} names the two routes to more: further repeated endpoints, or a
cohort attribute table.
\end{assumpbox}

The count is worth stating plainly. The model has $65$ values of $\delta$, $49$
of $\kappa$ and $29$ of $\eta$, against $A=118$ anchor rows. Sixteen observables
have a single anchor row each, because their source reports a centre only; for
those, $\delta_m$ absorbs the residual exactly.

Therefore $\mu$ is not identified by any single observable. It is identified by
the \emph{agreement} of residuals across observables, through the shrinkage of
$\delta$ toward zero. The same holds for $\omega$ through $\kappa$. The scales
$\tau_\delta$ and $\tau_\kappa$ do the work, and they are estimated from how
much the corrections disagree with each other.

This is a partial-pooling argument and it is sound. But the model does not have
$118$ independent measurements of the population, and no report should suggest
that it does.

\section{How to emit the population}
\label{sec:emit}

The fit gives a posterior over $\phi$. A virtual population is a set of patients.
The rule that turns one into the other is part of the model, and getting it wrong
throws away the distinction that section \ref{sec:centre} was careful to make.

The object that carries both quantities is the posterior predictive over
patients:
\begin{equation}
p(\theta_{\text{new}}\mid \text{data})
\;=\;\int F(\theta\mid\phi)\;p(\phi\mid\text{data})\;d\phi .
\label{eq:predictive}
\end{equation}
Its spread holds the variability between patients \emph{and} our uncertainty
about the population. Do not collapse the two by accident. We emit three
products, and each one carries its own label.

\textbf{The reference population.} Draw at one $\hat\phi$: the posterior mean of
$\mu$ and the posterior median of $\omega$. Its spread is between-patient
variability only. Use it when one population must be shipped, and say on the
label that it carries no uncertainty about $\phi$.

\textbf{The ensemble.} Keep $S$ populations, one per posterior draw of $\phi$.
Run the downstream study once per population. Report the spread of the
\emph{results} across populations. This keeps the two variances separate in the
answer, and it is the default deliverable.

\textbf{The pooled predictive.} Equation \eqref{eq:predictive} itself: draw
$\phi^{(s)}$, then draw patients from $F(\cdot\mid\phi^{(s)})$, then pool. It is
wider than the reference population. Use it only when the decision needs one
population that is right on average.

Four rules apply to all three.

\begin{enumerate}
\item \textbf{Draw fresh $z$.} The $N$ frozen patients are a device that makes
      the gradient work. They are not a population. Draw new $z$ for the
      deliverable, at whatever size the deliverable needs.
\item \textbf{Carry the corrections forward.} For an observable that the fit
      used, apply $\delta_m$ and $\kappa_m$ to the prediction. For a \emph{new}
      observable, draw $\delta_{\text{new}}\sim\Ncal(0,\tau_\delta^2)$ and
      $\log\kappa_{\text{new}}\sim\Ncal(\log\kappa_0,\tau_\kappa^2)$ from the
      fitted hyperparameters. This is what the hierarchical form is for. A
      prediction that drops the corrections claims that the model is correct,
      after a fit that measured how wrong it is.
\item \textbf{Filter for viability, and report the rate.} Score each drawn
      $\theta$ with the restriction classifier
      (\texttt{inference/restriction.py}), which predicts whether the simulator
      returns a usable output. Run the survivors through the full simulator, not
      the emulator. Report the rejection rate. A high rate means that $F$ puts
      mass outside the feasible set of the model. That is a fault in the fit and
      not a step in post-processing.
\item \textbf{Keep prevalence weighting for one job only.} If a downstream user
      needs the population to match a reported marginal exactly, reweight the
      emitted patients (\texttt{vpop/weighting.py}, Allen et al. 2016). Report
      that as an adjustment on top of the fit. It is never the fit.
\end{enumerate}

\section{Checks}
\label{sec:checks}

No number in this model is reportable until the checks below have run. They run
in this order. Each one gates the next, because a later verdict means nothing
while an earlier one fails.

\begin{enumerate}
\item \textbf{Simulation-based calibration on $\phi$.} Draw $\phi^*$ from the
      prior, simulate anchor rows through \eqref{eq:obs}, refit, and check that
      $\phi^*$ lands uniformly inside its own posterior. This tests the sampler
      and the likelihood code. It never sees the real data. If it fails, nothing
      below is readable.
\item \textbf{The width split.} Generate synthetic data with a known $\omega$ and
      a known uniform narrowing of the simulator. Fit. Confirm that $s$ takes the
      narrowing and $u$ takes the pattern. Report the recovered
      $(s,\log\kappa_0)$ pair against the truth. Report $\lambda_m$ for all $M_w$
      observables beside it.
\item \textbf{The frozen cloud.} Run the whole fit three times with three
      different frozen clouds $z_{1:N}$. Compare the posteriors. The difference
      is the bias that $\Sigma_{\text{MC}}$ hedges and does not remove.
\item \textbf{The plug-in for $\Sigma$.} Change $\phi_0$ by a factor of two and
      refit. Section \ref{sec:sigma} records the result of this test on the
      earlier form of $\Sigma$. Repeat it on the bootstrap form.
\item \textbf{The three CD8 cohorts.} Their three medians differ by a study
      effect plus a sampling error. Check that the disagreement matches what the
      bootstrap $\Sigma_{\text{samp}}$ predicts. This is a direct test of section
      \ref{sec:sigma}, and it is the only one available. Report the posterior
      for $\tau_\eta$ that these two contrasts produce, and compare it against
      its prior.
\item \textbf{Standardised residuals.} Compute $(y-\widetilde
      Q)_r/\sqrt{\Sigma_{rr}}$ for all $A$ rows. Look for structure by cohort, by
      observable class and by $n_c$. Structure by cohort means that the cohort
      terms of section \ref{sec:cohort} are doing the wrong job.
\item \textbf{Leave one cohort out.} Refit without cohort $c$, then predict its
      rows and check coverage. Repeat over cohorts. This is the only external
      check in the list, and it is the one that tests whether $\tau_\delta$ and
      $\tau_\kappa$ are honest.
\item \textbf{The emulator, where the fit reads it.} Section \ref{sec:emulator}.
      Draw from $F(\theta\mid\hat\phi)$, run the full simulator, and compare
      against the emulator. Report the result beside the design residual.
\item \textbf{The variance budget.} For every parameter, report which of three
      things produced its width: the data moved it, the prior returned it, or the
      pool could not represent it. Quote no per-parameter width without this
      table.
\item \textbf{Sensitivity to $R$.} Refit under the four candidates of section
      \ref{sec:R}, including the identity. Report how far $\hat\mu$ and
      $\hat\omega$ move. $R$ is assumed, so this number is part of the result.
\item \textbf{The viability rate.} Section \ref{sec:emit}, rule 3. Report the
      fraction of emitted patients that the simulator rejects.
\end{enumerate}

\section{What the fit estimates}

\begin{center}\small
\begin{tabularx}{\textwidth}{@{}llX@{}}
\toprule
Quantity & Status & Source \\
\midrule
$\mu$ & estimated, $P$ & Prior \eqref{eq:muprior} from the stage-1 mean and covariance. \\
$s$ & estimated, $1$ & The global width multiplier. Aliased with $\kappa_0$; see section \ref{sec:widthid}. \\
$u$ & estimated, $P{-}|\Mset|{-}1$ & The pattern of the widths. Sums to zero. \\
$\log\omega_\Mset$ & estimated, $9$ & Measured widths. Prior width from the source sample size. \\
$\eta,\ \delta$ & estimated & Centre corrections. Aliased; see section \ref{sec:alias}. \\
$\kappa,\ \kappa_0$ & estimated & Width corrections. See section \ref{sec:widthid}. \\
$\tau_\eta,\tau_\delta,\tau_\kappa$ & estimated & These carry the separation of the corrections. \\
$\varepsilon$ & fixed & Emulator residual, on the diagonal of $\Sigma$. Section \ref{sec:emulator}. \\
$R=\Corr(\Omega)$ & fixed & Section \ref{sec:R}. An input, not a result. \\
$\omega_0$ & fixed & Section \ref{sec:omega}. \\
$\mathcal{E}_c$ & fixed & Eligibility rules. Section \ref{sec:cohort}. \\
$\Sigma$ & fixed & Section \ref{sec:sigma}, at a plug-in $\phi_0$. \\
$z_{1:N}$ & fixed & Drawn one time. \\
\bottomrule
\end{tabularx}
\end{center}

The fit estimates about $690$ quantities against $A=118$ anchor rows.
At most $49$ functions of $\Omega$ are identified by the data, so the prior
carries most of $\omega$, and the shrinkage of the corrections carries the rest.
A count of free parameters is not a warning by itself in a hierarchical model.
It is a warning here, because section \ref{sec:alias} shows that the $118$ rows
are not $118$ independent measurements.

\section{Assumptions}
\label{sec:assump}

\begin{enumerate}[label=(\roman*)]
\item \textbf{We assume $R$ rather than estimating it.} The correlations of the
population are an input \eqref{eq:R}, and no part of them is identified by these
data. The current choice is the stage-1 correlation matrix, which is epistemic:
it shows how our uncertainties move together, not how patients differ. It is
also close to the identity. Section \ref{sec:R} lists better candidates, and
check 9 measures the sensitivity. That check has not yet been run.

\item \textbf{No data gives the correlations between patients.} Papers report
each observable on its own, so the cross terms $J_j\Omega J_l\Tr$ are never
observed. To measure them, we need two or more observables in the same
individuals. Digitised per-patient figures could supply this, but only if the
extraction keeps the identity of each patient across panels.

\item \textbf{The spread is mostly prior, not data.} At most $49$ functions of
$\Omega$ are identified against $271$ free widths. Also, $262$ of the $271$
values of $\omega_0$ come from a class default rather than a measurement. The
width of the reported population is therefore in large part a statement of
assumption. This is why $\omega_0$ must come from external evidence, and why
$\tau_\omega$ should not be widened for convenience.

\item \textbf{The global width and the global model error share one direction.}
Equation \eqref{eq:ridge} shows that $s$ and $\log\kappa_0$ enter every width row
through the sum $(1-\lambda_m)s+\log\kappa_0$. If $\lambda_m$ is near zero for
every observable, the prior width $\tau_s$ decides the split and the data do
not, because $s$ and $\log\kappa_0$ share it. Report $\lambda_m$ and the joint posterior of the pair. Check 2
tests the split on synthetic data, and it has not yet been run.

\item \textbf{The centre corrections are aliased and the priors separate them.}
Section \ref{sec:alias} shows that $\eta_c$ and $\delta_m$ are exactly aliased
for $64$ of the $65$ observables, because those observables appear in one cohort
each. The count of anchor rows overstates the information by a wide margin. Any
statement about $\mu$ or $\omega$ depends on $\tau_\delta$ and $\tau_\kappa$
remaining small, which is an assumption about how consistent the sources are.

The CD8 split of section \ref{sec:repeat} is the one exception, and it changes
the status of $\tau_\eta$ from asserted to measured. It rests on two contrasts.
Treat it as one weak measurement and not as a solution.

\item \textbf{Cohort differences are assumed to be assay effects.} $\eta_c$
shifts every readout of a study by a common amount, which is a statement about
calibration. A cohort that recruited different patients would instead shift the
parameters. For a cohort with one observable the two predict the same number, and
most cohorts have about two, so these data cannot separate them. The model takes
the assay reading by construction, and every statement about $\mu$ inherits that
choice. Section \ref{sec:cohort} explains why we removed the parameter-level term
rather than carrying both.

\item \textbf{The emulator error term is specified and not yet populated.} The
slot $\varepsilon_m$ enters $\Sigma$ through \eqref{eq:Sigma}. In the current
pipeline nothing writes to it, so it is zero. The emulator's held-out residual is
not negligible: relative to the within-cohort spread it lies between $0.095$ and
$0.254$. Until $\varepsilon$ is wired, every target is trusted more than it
should be, and the highest-$n$ targets most of all. Section \ref{sec:emulator}
also requires that this residual be re-measured on draws from the fitted
population, and that has not been done.

\item \textbf{The measured spreads are biased in both directions.} A reported
across-patient spread also contains assay noise and biopsy-site variation, which
makes it too large. A spread measured across animals of one genetic background
is smaller than a spread across human patients, which makes it too small. We
correct neither effect in \eqref{eq:shrink}.

\item \textbf{Parts of this note describe the model and not the code.} The
document is the specification. At the time of writing,
\texttt{vpop/population\_fit.py} still uses the basis form
$\vartheta=\mu+W(\sigma_u\odot z)$ and has no $\kappa$, no $s$ and no $u$, no
eligibility rule, and an asymptotic $\Sigma$. The items to build are: the
$(s,u,\Mset)$ split of section \ref{sec:spread}; the two-row anchors of
\eqref{eq:Q}; the additive $\log\kappa$ of \eqref{eq:corrwid}; the bootstrap
$\Sigma_{\text{samp}}$ and the measured $\Sigma_{\text{MC}}$ of section
\ref{sec:sigma}; the cohort terms of section \ref{sec:cohort}; the stage-1
covariance in \eqref{eq:muprior}; and the emission rules of section
\ref{sec:emit}. The CD8 split of section \ref{sec:repeat} is also outstanding,
and every count in this note that involves $C$ or $A$ assumes it. Do that one
first. It is extraction work, it needs no new code path beyond a cohort index,
and it is the only change here that adds information rather than structure.
\end{enumerate}

\section{What data to gather next}
\label{sec:data}

Section \ref{sec:assump} lists nine assumptions. Most of them are not defects of
the method. They are statements that the available data cannot answer a
question. This section says which data would answer which one, and in what
order to go and get it.

\begin{notebox}
\textbf{The rule that sets the order.} A new anchor row is worth what it shares
with a row that you already hold.

A row that shares an \emph{endpoint} with another study identifies $\eta$. A row
that shares a \emph{cohort} with another observable separates an assay effect
from a recruitment effect, and helps $\mu$. A row that shares a \emph{patient}
with another observable identifies $R$. A row that shares nothing arrives with
its own $\delta_m$ and its own $\kappa_m$, so it contributes only through the
shrinkage of the corrections.

This rule reverses the usual instinct. One more observable, in one more cohort,
at the same three quantile levels, is the least valuable thing that you can add.
It is also the easiest thing to add.
\end{notebox}

\subsection{Price the options before you buy them}

Three checks in section \ref{sec:checks} cost almost nothing, and they set the
ranking below. Run them first.

\textbf{Check 9, the sensitivity to $R$.} Refit under the identity and under the
stage-1 matrix. If $\hat\mu$ and $\hat\omega$ barely move, then item 1 below is
worth less than it claims. If they move, item 1 is the whole question. You
cannot know which without the run.

\textbf{The $\lambda_m$ report of section \ref{sec:widthid}.} It says whether
measured parameter spreads can break the ridge between $s$ and $\log\kappa_0$,
and which observables would do it.

\textbf{Check 8 and the emulator resolution of section \ref{sec:emulator}.} A
parameter that the emulator cannot resolve gains nothing from a better prior.
Do not buy data for those parameters.

\subsection{The ranking}

\begin{enumerate}
\item \textbf{Two or more observables in the same patients.} This is the only
source of $R$, and it is the only item that attacks an assumption that section
\ref{sec:whatdataconstrain} calls unreachable from summary data. It does four
jobs at once. It gives the cross terms $J_j\Omega J_l\Tr$. It gives a full
marginal instead of three quantiles. It puts several observables in one cohort.
It removes the plug-in from $\Sigma$.

The cheap form of this item exists today. Papers in this area print per-patient
dot plots. Digitise them, and keep the identity of each patient across panels.
Twenty patients with five observables each would change this model more than
twenty new targets. The extraction discipline is the whole requirement. A panel
digitised without patient identity is worth a fraction of the same panel
digitised with it. Removes (i) and (ii).

\item \textbf{Repeated endpoints across studies.} The cheapest item here, and
the CD8 target of section \ref{sec:repeat} proves that the stock exists. For each
of the $65$ observables, search for a second publication that reports the same
quantity in a comparable cohort. Five repeated endpoints move section
\ref{sec:alias} from a prior structure to a measurement, and $\tau_\eta$ then
governs how the model splits every other residual between a study effect and a
model error. Removes (v).

\item \textbf{Across-patient spreads for parameters, ordered by sensitivity.}
Nine of $271$ parameters have a measured spread. Each one that you add does two
jobs: it improves $\omega_0$, and it enters $\Mset$, which is what makes
$\lambda_m$ vary and breaks the width ridge. The second job is why the early
returns are large.

Do not search at random. Rank the parameters by their contribution to the
predicted variance of the observables that have a width row. That ranking is
computable today from $J_m$, $R$ and $\omega_0$. Then find donor-to-donor or
patient-to-patient spreads for the top of the list, with sample sizes. The gate
of section \ref{sec:omega} still applies: biological units only, and never
replicates. Removes part of (iii) and part of (iv).

\item \textbf{More quantile levels from the papers that you already cite.} With
three levels, $\kappa_m$ copies any width change exactly. With a reported range,
a tenth and a ninetieth centile, or a mean printed beside a median, it cannot,
because $\omega$ changes the shape of the pushforward and $\kappa$ only stretches
it. This is a re-read of sources that are already in the set. It is the cheapest
attack on the largest remaining ambiguity. Removes part of (iv).

\item \textbf{Cohort attributes and eligibility criteria.} For all $29$ cohorts:
stage, line of therapy, treatment status, sample type, assay platform, region.
This turns $29$ free values of $\eta_c$ into $Z_c\beta$, whose coefficients pool
across cohorts, and it is also what a covariate model on the parameters would
need. A quantitative eligibility rule is better still, because it removes a
parameter instead of constraining one. This is data entry against sources that we
hold. Removes part of (v) and part of (vi).

The attribute table also does one thing that no other item does. An attribute
that is plainly instrumental, such as the assay platform, and an attribute that
is plainly clinical, such as the stage, load on the same $\eta_c$ today. Split
them and the assay reading of assumption (vi) becomes testable rather than
asserted.

\item \textbf{Repeat measurements on the same patient.} Assumption (viii) says
that every reported across-patient spread also holds assay noise, and that we
correct none of it. So every $\omega_0$ that L4 built is too large by an unknown
amount. A technical replicate, or a paired sample from one patient, measures that
noise directly and lets us subtract it. This is the only item in the list that
removes a known bias rather than a known ambiguity. Removes (viii).
\end{enumerate}

\subsection{Summary}

\begin{center}\small
\begin{tabularx}{\textwidth}{@{}Xll@{}}
\toprule
Data & Removes & Cost \\
\midrule
Per-patient values, several observables, identity kept & (i), (ii) & digitisation \\
The same endpoint in a second study & (v) & literature search \\
Parameter spreads with $n$, ranked by sensitivity & (iii), (iv) & literature search \\
Extra quantile levels from current sources & (iv) & re-read \\
Cohort attributes and eligibility rules & (v), (vi) & data entry \\
Replicate or paired samples & (viii) & new measurement \\
\bottomrule
\end{tabularx}
\end{center}

\subsection{If we can commission one measurement}

One cohort. One platform. Five to ten of the model's observables, measured in the
same patients, with the per-patient values published. That is item 1, item 5 and
part of item 2 in a single study. Nothing else on the list gives three at once.

\subsection{What not to buy}

A new observable, in a new cohort, at three quantile levels. It arrives with a
new $\eta_c$, a new $\delta_m$ and a new $\kappa_m$. The model gains one row and
three parameters, and the information reaches $\mu$ and $\omega$ only through
$\tau_\delta$ and $\tau_\kappa$.

The same observable, placed in a cohort that is already in the set, is worth
several times more. It adds no $\eta$, and it gives that cohort a second
observable, which is what assumption (vi) needs before an assay effect and a
recruitment effect can be told apart. When there is a choice of source, that is
the choice.

\subsection{Make this list computable}

\texttt{inference/obed.py} ranks candidate measurements by expected information
gain. Point it at $(\mu,\ s,\ \log\kappa_0,\ \log\omega_\Mset)$ rather than at
the full $\phi$, and this section becomes a number instead of a judgement. The
ranking above is what we expect it to return. The run is how we find out where
the expectation is wrong.

\end{document}
